% =============================================================================
% Fig 1: OrbitVLIF system overview
%   (a) VLIFNet architecture pipeline (single satellite)
%   (b) Decentralised federated topology over 5×10 LEO constellation
%
% Compile:
%   pdflatex fig1.tex
%
% Style: IEEE 2-column conference, vector-only, Wong-2011 color-blind palette,
%        STIX math fonts, sans-serif body.  Designed for ~3.5"×2.4" rendering.
% =============================================================================
\documentclass[border=4pt, tikz]{standalone}

\usepackage[T1]{fontenc}
\usepackage{stix}                     % math + body in STIX (Times-like)
\usepackage{amsmath, amssymb, bm}
\usepackage{tikz}
\usetikzlibrary{
    positioning, arrows.meta, shapes.geometric, shapes.arrows,
    fit, backgrounds, calc, decorations.pathreplacing, decorations.markings,
    patterns, patterns.meta
}

% ---- Color-blind safe palette (Wong 2011) -----------------------------------
\definecolor{cBlue}  {HTML}{0173B2}     % VLIFNet / SNN
\definecolor{cOrange}{HTML}{DE8F05}     % ANN / Plain
\definecolor{cGreen} {HTML}{029E73}     % FSTA / accent 1
\definecolor{cRed}   {HTML}{D55E00}     % attention / accent 2
\definecolor{cPurple}{HTML}{CC78BC}     % MultiSpike / accent 3
\definecolor{cGray}  {HTML}{4D4D4D}     % text / outline
\definecolor{cLight} {HTML}{E8E8E8}     % subtle background

% ---- Block style helpers ----------------------------------------------------
\tikzset{
  font=\sffamily\footnotesize,
  >={Stealth[length=4pt, width=3pt]},
  every node/.style={inner sep=2pt, outer sep=0pt, align=center},
  block/.style={
    rectangle, rounded corners=1.5pt, draw=cGray, line width=0.4pt,
    minimum width=14mm, minimum height=5mm,
  },
  encblk/.style={block, fill=cBlue!18, draw=cBlue!70},
  decblk/.style={block, fill=cOrange!18, draw=cOrange!70},
  attblk/.style={block, fill=cGreen!18, draw=cGreen!70, font=\sffamily\scriptsize},
  ioblk/.style ={block, fill=cLight, draw=cGray, minimum width=10mm},
  fusblk/.style={block, fill=cPurple!18, draw=cPurple!70, font=\sffamily\scriptsize,
                 minimum width=8mm, minimum height=4mm},
  arrow/.style={->, line width=0.6pt, draw=cGray, shorten >=1pt, shorten <=1pt},
  skip/.style ={->, line width=0.5pt, draw=cPurple!70, dashed,
                shorten >=1pt, shorten <=1pt},
  label/.style={font=\sffamily\scriptsize, text=cGray},
  panellbl/.style={font=\sffamily\bfseries\small, text=black},
  satnode/.style={
    circle, draw=cGray, line width=0.3pt, fill=white,
    minimum size=2.6mm, inner sep=0pt, font=\sffamily\tiny,
  },
  satnodeS/.style={
    circle, draw=cBlue, line width=0.45pt, fill=cBlue!25,
    minimum size=2.8mm, inner sep=0pt,
  },
  planebox/.style={
    rectangle, rounded corners=2pt, draw=cGray!50, line width=0.3pt,
    fill=cLight!50, inner sep=2.5pt,
  },
  ringedge/.style={-, line width=0.3pt, draw=cGreen!75, opacity=0.85},
  mdstedge/.style={-, line width=0.5pt, draw=cRed!85,
                   decoration={markings, mark=at position 0.55 with {\arrow{Stealth[length=3pt]}}},
                   postaction={decorate}},
}

\begin{document}
\begin{tikzpicture}

% =============================================================================
% Panel (a): VLIFNet architecture pipeline
% =============================================================================

% ---- Input / output -----------------------------------------------------------
\node[ioblk] (in)  at (0, 0)         {Cloudy\\$\mathbf{x}$};
\node[ioblk] (out) at (12.6, 0)      {Restored\\$\hat{\mathbf{y}}$};

% ---- Patch embed ---------------------------------------------------------------
\node[block, fill=cGray!10, minimum width=10mm] (pe) [right=3mm of in]  {Patch\\Embed};

% ---- Encoder (3 levels) -------------------------------------------------------
\node[encblk, minimum width=12mm] (e1) [right=3mm of pe] {SRB$\times2$\\$\mathit{C}=24$};
\node[encblk, minimum width=12mm] (e2) [right=3mm of e1] {SRB$\times2$\\$\mathit{C}=48$};
\node[encblk, minimum width=12mm] (e3) [right=3mm of e2] {SRB$\times4$\\$\mathit{C}=96$};

% Down arrows + "/2" labels between encoder levels
\draw[arrow] (pe) -- (e1);
\draw[arrow] (e1) -- node[label, above]{$\downarrow2$} (e2);
\draw[arrow] (e2) -- node[label, above]{$\downarrow2$} (e3);

% ---- Bottleneck (level-1 freq block) ------------------------------------------
\node[fusblk, fill=cRed!18, draw=cRed!75, minimum width=14mm, minimum height=6mm,
      font=\sffamily\scriptsize]
      (lvl1) [right=3mm of e3] {$\mathcal{L}_1$ Block\\(MS-4 + MDA)};
\draw[arrow] (e3) -- (lvl1);

% ---- Decoder (3 levels, mirrored, drawn below encoder) ------------------------
\node[decblk, minimum width=12mm] (d3) [below=8mm of e3]   {SRB$\times2$};
\node[decblk, minimum width=12mm] (d2) [left=3mm of d3]    {SRB$\times2$};
\node[decblk, minimum width=12mm] (d1) [left=3mm of d2]    {SRB$\times2$};

% Bottleneck connects to decoder L3
\draw[arrow] (lvl1.south) |- (d3.east);
% Up arrows
\draw[arrow] (d3) -- node[label, above]{$\uparrow2$} (d2);
\draw[arrow] (d2) -- node[label, above]{$\uparrow2$} (d1);

% Output head
\node[block, fill=cGray!10, minimum width=10mm] (head) [left=3mm of d1] {Output\\Head};
\draw[arrow] (d1) -- (head);
\draw[arrow] (head.west) -- ++(-3mm,0) |- (out.west);

% Residual connection from input to output
\draw[arrow, draw=cGray, dashed]
  (in.north) -- ++(0,4mm) -- node[label, above]{global residual} ($(out.north)+(0,4mm)$) -- (out.north);

% ---- Skip connections (encoder → decoder via GatedSkipFusion) ----------------
% e1 -> d1 (long skip)
\draw[skip]
  (e1.south) -- ++(0,-3mm) -| (d1.north);
% e2 -> d2
\draw[skip]
  (e2.south) -- ++(0,-2mm) -| (d2.north);
% e3 -> d3
\draw[skip]
  (e3.south) -- ++(0,-1mm) -| (d3.north);

% Label for skip path style
\node[label, text=cPurple, font=\sffamily\scriptsize]
  at ($(e2.south)+(0,-4mm)$) {GatedSkipFusion};

% ---- FSTA + MultiSpike-4 callout boxes (legend overlays) ----------------------
\begin{scope}[every node/.style={font=\sffamily\scriptsize, align=center}]
  \node[attblk, minimum width=22mm, minimum height=6mm]
        (fsta) at ($(e3.north)+(0,9mm)$) {FSTA\\(Temporal Amp + DCT Spatial)};
  \draw[arrow, draw=cGreen!70, line width=0.4pt]
        (fsta.south) -- ($(e3.north)+(0,1pt)$);

  \node[block, fill=cPurple!18, draw=cPurple!70, minimum width=22mm, minimum height=5mm,
        font=\sffamily\scriptsize]
        (ms4) at ($(lvl1.north)+(0,9mm)$) {MultiSpike-4 quantizer\\$\{0,\tfrac14,\tfrac12,\tfrac34,1\}$};
  \draw[arrow, draw=cPurple!70, line width=0.4pt]
        (ms4.south) -- ($(lvl1.north)+(0,1pt)$);
\end{scope}

% ---- Panel (a) label ----------------------------------------------------------
\node[panellbl] at (-0.3, 1.8) {(a) VLIFNet architecture (single satellite)};

% =============================================================================
% Panel (b): Federated topology (orbit constellation, 5 planes × 10 sats)
% =============================================================================

% Place panel (b) below panel (a) with vertical gap
\begin{scope}[shift={(0, -5.6)}]
  \node[panellbl] at (-0.3, 2.4)
        {(b) Decentralised satellite federation: $5$ orbital planes $\times\,10$ satellites};

  % Define 5 planes laid horizontally
  \def\nsat{10}
  \def\nplane{5}

  \foreach \p in {0,1,2,3,4} {
    % Plane y-position
    \pgfmathsetmacro{\px}{2.4*\p + 0.6}

    % Plane bounding box (drawn first as background)
    \begin{pgfonlayer}{background}
      \node[planebox, minimum width=22mm, minimum height=22mm]
            (plane\p) at (\px, 0) {};
    \end{pgfonlayer}

    % Label above each plane
    \node[label, text=cGray, anchor=south]
          at ($(plane\p.north)+(0,1pt)$) {Plane $\mathcal{P}_{\p}$};

    % Place 10 satellites on a circle inside this plane
    \foreach \i in {0,...,9} {
      \pgfmathsetmacro{\ang}{(\i / \nsat) * 360}
      \coordinate (P\p-S\i) at ($(plane\p.center) + (\ang:7.5mm)$);
      \ifnum\i=0
        \node[satnodeS] at (P\p-S\i) {};
      \else
        \node[satnode]  at (P\p-S\i) {};
      \fi
    }

    % Intra-plane gossip ring (connect adjacent satellites in cycle)
    \foreach \i [evaluate=\i as \j using int(mod(\i+1,10))] in {0,...,9} {
      \draw[ringedge] (P\p-S\i) -- (P\p-S\j);
    }
  }

  % MDST inter-plane edges (selected, not full mesh — represents minimum spanning tree)
  % Plane 0 ↔ 1, 1 ↔ 2, 2 ↔ 3, 3 ↔ 4 (root chain)
  \foreach \a/\b in {0/1, 1/2, 2/3, 3/4} {
    \draw[mdstedge] (P\a-S0) to[bend left=20] (P\b-S0);
  }
  % One additional inter-plane shortcut (typical MDST has > N-1 redundancy)
  \draw[mdstedge, draw=cRed!50, dashed, line width=0.4pt]
        (P0-S0) to[bend left=35] (P4-S0);

  % Legend (right of last plane)
  \begin{scope}[shift={(13.0, 0.4)}]
    \node[satnode]  (lg1) at (0,1.0) {};
    \node[label, anchor=west] at ($(lg1.east)+(2pt,0)$) {Sat. (worker)};
    \node[satnodeS] (lg2) at (0,0.6) {};
    \node[label, anchor=west] at ($(lg2.east)+(2pt,0)$) {Plane head};
    \draw[ringedge, line width=0.7pt] (-1.5mm, 0.2) -- (1.5mm, 0.2);
    \node[label, anchor=west] at (3pt, 0.2) {Intra-plane gossip};
    \draw[mdstedge, line width=0.7pt] (-1.5mm,-0.2) -- (1.5mm,-0.2);
    \node[label, anchor=west] at (3pt, -0.2) {Inter-plane MDST};
  \end{scope}

  % FedBN annotation under the planes
  \node[label, text=cGray, font=\sffamily\scriptsize, align=center]
        at (6, -1.8)
        {FedBN $+$ Gossip $+$ Dirichlet$(\alpha\!=\!0.1)$ partition over cloud-type label};

\end{scope}

\end{tikzpicture}
\end{document}
